\documentclass{article}
\usepackage[T1]{fontenc}
\usepackage{lmodern}
\usepackage{graphicx}
\usepackage{amsmath,amssymb}
\usepackage[a4paper,margin=2cm]{geometry}
\usepackage[absolute,overlay]{textpos}
\setlength{\TPHorizModule}{1mm}
\setlength{\TPVertModule}{1mm}
\usepackage[indonesian]{babel}
\usepackage{hyperref}
\setlength{\emergencystretch}{2em}
% unit: Halaman 15

\begin{document}

% === Halaman 15 (python/native) ===

15

• Contoh 3: Misalkan kriptanalis mencoba kunci  LMCCAWAAZD

   untuk mendekripsi cipherteks  HOJKOREGHP

 Plainteks yang dihasilkan: SALMONEGGS    (telur ikan Salmon)

   Bila ia mencoba kunci: ZDVUZOEYEO

 Plainteks yang dihasilkan: GREENFIELD    (lapangan hijau)

 Kriptanalis: ???????   Kunci mana yang benar? (bingung sendiri ☺ )

• Contoh ini menunjukkan bahwa untuk sembarang plainteks dan 

cipherteks hanya ada satu kunci yang memetakannya satu sama lain.

\end{document}
